% Pursuing Stacks (cotes 134-2 to 134-8): archive cote 134-2, site page 38 (PDF page 39).
% appendix (letters to Breen). Draft compared: ps1.tex:1174-1193.
% AI-transcribed by Claude Opus 5.5 (claude-opus-5-5, adaptive thinking, effort high), Message Batches API, batch ps1 (msgbatch_01Rv5hCaNTtUt4JHx1yqoHPL), 2026-11-12; not yet checked by eye.
\scanpage{134-2/38}
\mspage{\hand{23 bis}}
% NOTE: the typed text is on a slip pasted onto another typed page; text of the underlying page shows through faintly and is not transcribed

\hfill Villecun 5.2.1975

\margin{--} Cher Breen

\quad ....

(App.1) \quad ... Pour tout le reste de ta lettre, \ elle mériterait un lecteur
\margin{--} plus averti, \ aussi, \ pour qu'elle ne soit pas entièrement perdue au monde,
je vais l'envoyer à Illusie ! \ J'ai néanmoins
constaté\hand{,} avec intérêt\hand{,} ton intérêt à demi \unsure{rèfoulé|refoulé} pour des 2-catégories de
Picard, n-catégories et autre faune de ce genre, et ton espoir que je te
prouverai peut-être que ces animaux sont tout à fait indispensables pour
faire des maths sérieuses \ dans telle ou telle circonstance. J'ai bien \hand{peur}
que cet espoir ne soit déçu, je crois que jusqu'à maintenant on a toujours
pu s'en tirer en éludant de tels objets et l'engrenage dans lequel ils pour
-raient nous entrainer. Est-ce nécessairement une raison pour continuer à
les éluder ? Les situations où on a l'impression"d'éluder" en effet me sem-%
blent en tous cas devenir toujours plus nombreuses - \ et si on s'abstenait
de tirer une situation complexe et chargée de mystère au clair, chaque
fois qu'on ne serait pas \uline{forcé} de le faire pour des raisons techniques pro-%
venant de la math déjà faite, - il y aurait sans doute beaucoup de parties
des maths aujourd'hui reputées\hand{|}"sérieuses" qui n'auraient jamais \inserted{été} développées
(Il n'est pas dit non plus que\hand{|}le monde s'en trouverait plus mal ...). Ton
commentaire (que j'ai également entendu chez Deligne) \ que la classificati-%
on d'objets géométriques relativement merdiques se réduit finalement à des
invariants cohomologiques essentiellement "bien connus" et relativement
simples n'est pas non plus convainquant; n'est ce pas négliger la différen-%
ce entre la \uline{compréhension d'un objet géométrique}, et la détermination de
sa "classe à isomorphisme (ou équivalence...) près" ?

\margin{38}
